% Source-independent fallback. This does not validate either unrecovered note.
% Requires amsmath, amssymb, amsthm and theorem/corollary/proposition environments.
\section{An exact modulus criterion by coefficient moments}
\label{v2:sec:moments}

The classification of the mates of a specified set and the classification of
moduli have different quantifiers.  The following finite criterion addresses
the latter directly.  Its proof is independent of the sign-graph theorem.  It
is an exact enumeration and aggregation procedure, with exponential worst-case
cost; no claim of a short arithmetic characterization is made.

Write $\mathcal D_N=\{d:d\mid N\}$ and let $\mathcal J_N$ be the family of
downward-closed subsets of $\mathcal D_N$, including the empty subset.  For
$I\in\mathcal J_N$, put
\[
 A_I=\{x\in\mathbb Z_N:\operatorname{ord}(x)\in I\}.
\]
Thus $I\mapsto A_I$ is a bijection from $\mathcal J_N$ to subgroup unions.
In this section every divisor indexing a physical shell denotes an
\emph{element order}, not a subgroup index.  Define the integer Ramanujan sum
\[
 c_d(k)=\sum_{h\mid\gcd(d,k)}h\mu(d/h).
\]
The Fourier value at character conductor $c\mid N$ and its squared signature
are
\begin{equation}\label{eq:moment-signature}
 F_I(c)=\sum_{d\in I}c_d(N/c),\qquad
 \sigma(I)=\bigl(F_I(c)^2\bigr)_{c\mid N}.
\end{equation}
The character sum over the subgroup of order $d$ is $d$ when $d$ divides the frequency and zero otherwise. Subtracting the contributions of proper order shells, or equivalently applying divisor M\"obius inversion, gives the displayed formula for $c_d(k)$, where $\mu$ is the M\"obius function. Summing the shell contributions proves the first identity.
Every $A_I$ is inversion symmetric, so Fourier inversion gives
\begin{equation}\label{eq:moment-signature-equivalence}
 C_{A_I}=C_{A_J}\quad\Longleftrightarrow\quad\sigma(I)=\sigma(J).
\end{equation}
This compares all conductors, including mixed conductors.

For one indeterminate $z_c$ per $c\mid N$, use
$z^v=\prod_{c\mid N}z_c^{v_c}$.  For a physical order $d\mid N$, set
\begin{equation}\label{eq:cut-polynomials}
 P_d^+(z)=\sum_{\substack{I\in\mathcal J_N\\d\in I}}z^{\sigma(I)},
 \qquad
 P_d^-(z)=\sum_{\substack{I\in\mathcal J_N\\d\notin I}}z^{\sigma(I)}.
\end{equation}
Their coefficients are nonnegative integers.  For a finite Laurent polynomial
$P=\sum_v a_vz^v$, write $\|P\|_{\mathrm{coef}}^2=\sum_v a_v^2$ and
let $[z^0]P$ denote its constant coefficient.

\begin{theorem}[Cut coefficient criterion]\label{thm:cut-coefficient}
For every positive integer $N$ and divisor $d\mid N$, the integer
\begin{align}
 K_d(N)&=[z^0]\bigl(P_d^+(z)P_d^-(z^{-1})\bigr)\label{eq:cut-moment}\\
 &=\frac14\left(
 \|P_d^++P_d^-\|_{\mathrm{coef}}^2-
 \|P_d^+-P_d^-\|_{\mathrm{coef}}^2\right)\notag
\end{align}
is nonnegative.  It is positive if and only if there are homometric subgroup
unions $A_I,A_J$ with $d\in I$ and $d\notin J$.  Consequently, $N$ is rigid
among subgroup unions if and only if $K_d(N)=0$ for every $d\mid N$.
A positive moment yields an explicit pair of distinct homometric subgroup
unions.
\end{theorem}
\begin{proof}
Let $a_v=[z^v]P_d^+$ and $b_v=[z^v]P_d^-$.  The constant coefficient in
\eqref{eq:cut-moment} is $\sum_v a_vb_v$.  Each summand is nonnegative,
and $(a_v+b_v)^2-(a_v-b_v)^2=4a_vb_v$ proves the second formula.
Thus $K_d(N)>0$ precisely when one signature $v$ occurs in both cut classes.
Choosing an ideal from each class gives two different ideals whose
autocorrelations agree by \eqref{eq:moment-signature-equivalence}.
Conversely, two different homometric ideals differ in membership at some
$d\mid N$, and, after interchanging them if necessary, contribute a positive
term to $K_d(N)$.  The constructions below retain witnesses for every
nonzero coefficient, making the choice effective.
\end{proof}

The cuts with zero moment are exactly those order-shell memberships that are determined uniformly across every possible autocorrelation at the given modulus. For a specified signature, the same coefficient test can instead be restricted to that signature.  The sum of
the $K_d(N)$ counts each unordered pair of distinct homometric ideals with
weight $|I\mathbin\triangle J|$.  Thus the test concerns actual binary,
divisor-closed sets throughout.  It does not replace those conditions by
linear or rational relaxations.

\subsection{A prime-factor transfer recurrence}

We specify an exact way to form \eqref{eq:cut-polynomials}.  This avoids
testing every pair of ideals for homometry, but it still represents the full
finite ideal family and can require exponential resources.

For $N>1$, choose a prime-power factorization $N=p^aM$ with
$a\ge1$ and $\gcd(p,M)=1$; $M=1$ is allowed.  Each ideal $I\in\mathcal J_N$
has unique sections
\[
 I_t=\{m\mid M:p^tm\in I\},\qquad 0\le t\le a,
\]
and downward closure is equivalent to
\begin{equation}\label{eq:ideal-flag}
 I_0\supseteq I_1\supseteq\cdots\supseteq I_a,
 \qquad I_t\in\mathcal J_M.
\end{equation}
For $0\le j,t\le a$, define
\[
 r_{j,t}=c_{p^t}(p^{a-j}).
\]
For $t=0$ this is $1$; for $t\ge1$ it is $\varphi(p^t)$ if
$j\le a-t$, is $-p^{t-1}$ if $j=a-t+1$, and is zero otherwise.
For $J\in\mathcal J_M$, let the vector $V_t(J)$, indexed by
$(j,c)\in\{0,\ldots,a\}\times\mathcal D_M$, be
\begin{equation}\label{eq:factor-contribution}
 V_t(J)_{j,c}=r_{j,t}\sum_{m\in J}c_m(M/c).
\end{equation}
Multiplicativity of Ramanujan sums, or a direct character calculation on
$\mathbb Z_{p^a}\times\mathbb Z_M$, gives
\begin{equation}\label{eq:factor-spectral-sum}
 F_I(p^jc)=\sum_{t=0}^a V_t(I_t)_{j,c}.
\end{equation}

Fix a cut $d=p^sm$, with $m\mid M$, and $\epsilon\in\{0,1\}$.
Put $b_t^{\epsilon}(J)=1$ except when $t=s$; at $t=s$ put
$b_s^{\epsilon}(J)=1$ if $\mathbf1_J(m)=\epsilon$ and zero otherwise.
Using Laurent monomials whose exponent vector is the \emph{unsquared}
Fourier data, define
\begin{align}
 Q_0^{\epsilon}(J;z)&=b_0^{\epsilon}(J)z^{V_0(J)},\label{eq:transfer-start}\\
 Q_t^{\epsilon}(J;z)&=b_t^{\epsilon}(J)z^{V_t(J)}
          \sum_{\substack{K\in\mathcal J_M\\K\supseteq J}}
                    Q_{t-1}^{\epsilon}(K;z),\qquad 1\le t\le a.
 \label{eq:transfer-step}
\end{align}
Let $\mathsf{Sq}$ be the additive map on finite Laurent polynomials defined
by $\mathsf{Sq}(z^v)=z^{(v_c^2)_c}$.  This map is not multiplicative: it
must be applied only after all section contributions have been added.

\begin{proposition}[Exact transfer and witness extraction]
\label{prop:cut-transfer}
With the foregoing definitions,
\[
 P_d^+=\mathsf{Sq}\left(\sum_{J\in\mathcal J_M}Q_a^1(J;z)\right),
 \qquad
 P_d^-=\mathsf{Sq}\left(\sum_{J\in\mathcal J_M}Q_a^0(J;z)\right).
\]
These formulas, followed by coefficient matching, decide rigidity of every
factored modulus $N$ and construct a witnessing pair whenever $N$ is
nonrigid.  If $L=\tau(N)$, a direct implementation has a finite bound
$2^{O(L)}\operatorname{poly}(L,\log N)$ in bit operations.  This is not a
polynomial bound in the binary input length $\log N$.
\end{proposition}
\begin{proof}
Inductively, $Q_t^{\epsilon}(J;z)$ is the sum over all flags
$I_0\supseteq\cdots\supseteq I_t=J$ satisfying the cut restriction if
$s\le t$, with monomial exponent $\sum_{u=0}^t V_u(I_u)$.
The claim for $t=0$ is \eqref{eq:transfer-start}.  Every flag of length
$t+1$ has exactly one penultimate section $K\supseteq J$, proving the
inductive step without omission or duplication.  At $t=a$, the flags
are exactly all ideals of $\mathcal D_N$ by \eqref{eq:ideal-flag}.
Equation \eqref{eq:factor-spectral-sum} identifies their exponents as
their complete Fourier vectors.  Applying $\mathsf{Sq}$ therefore gives
exactly the polynomials in \eqref{eq:cut-polynomials}.

All coefficients remain nonnegative in this recurrence.  One may retain
one generating flag for every nonzero coefficient.  When a squared
signature occurs in both cut classes, its two stored flags explicitly
give the required ideals and hence the sets $A_I,A_J$.  Alternatively,
retain the corresponding unsquared Fourier vectors and apply
\[
 \mathbf1_{A_I}(N/d)=\frac1N\sum_{c\mid N}F_I(c)c_c(N/d),\qquad d\mid N.
\]
The retained vectors were generated from flags of ideals, so this inverse
has binary entries and is downward closed; these properties are not being
inferred merely from equality of squared signatures.

For completeness, the state set $\mathcal J_M$ can be constructed by
testing the $2^{\tau(M)}$ subsets for downward closure.  There are at most
$2^{\tau(M)}$ states and $2^{2\tau(M)}$ transitions at any step, and
at most $2^L$ monomials in a state.  All exponent vectors have length
$L$.  Before squaring their coordinates have absolute value at most $N$,
because they are Fourier values of the union formed by the retained
sections; after squaring they are at most $N^2$.  Coefficients are bounded
by $2^L$.  Exact sorting and addition of these integer vectors and
coefficients, for at most $a+1\le L$ steps and $L$ cuts, gives the stated
exponential bound.  The case $N=1$ consists of the two ideals
$\varnothing$ and $\{1\}$ and is checked directly.
\end{proof}

\subsection{What the universal criterion does and does not classify}

Theorem~\ref{thm:cut-coefficient} has no restriction on the number or size
of the prime factors or on their exponents.  Its necessity, sufficiency,
and witness extraction are proved above.  It provides a finite answer
for each modulus once its factorization is supplied.  It does not bound
the amount of computation by a practical polynomial, and it does not
characterize the infinitely many rigid moduli by a concise set of
prime-factor conditions.  The latter remains a separate arithmetic
problem.  The transfer recurrence is a transparent universal baseline;
the useful structural reductions elsewhere in the paper can substantially
reduce the domain on which it would have to be run.

